% ID: SEL-D0074-OB-E01
% ESTADO: cotejada. Segunda pasada visual y borrador compilado revisados.
% Fuente: Selectividad/sel_2008_matematicas_aplicadas/Mat CCSS 2 - Examen y criterios de correción.pdf
% SHA-256: 6439a5956271f683dc45ba870695a7399db4412c9b0eaf1a35f908a3857ae648
% Página: 2; opción B; ejercicio original 1.

\textbf{(3 puntos)} Un nutricionista informa a un individuo que, en cualquier tratamiento que siga, no debe ingerir diariamente más de 240 mg de hierro ni más de 200 mg de vitamina B. Para ello están disponibles píldoras de dos marcas, P y Q. Cada píldora de la marca P contiene 40 mg de hierro y 10 mg de vitamina B, y cuesta 6 céntimos de euro; cada píldora de la marca Q contiene 10 mg de hierro y 20 mg de vitamina B, y cuesta 8 céntimos de euro.

Entre los distintos tratamientos, \textquestiondown{}cuál sería el de máximo coste diario?
